\chapter{Evaluation Methodology}
\label{chap:eval}

This chapter describes the dataset, evaluation metric, experimental baselines, and procedures used to assess the proposed method. The evaluation framework is designed to provide trustworthy performance estimates through local validation before any public leaderboard submission, ensuring that architectural decisions are grounded in evidence rather than overfitting to the public test set.

% =====================================================================
\section{Dataset}
\label{sec:dataset}
% =====================================================================

% =====================================================================
\subsection{Data Source and Split}
\label{sec:data_source}
% =====================================================================

All experiments use the official dataset from the ADIA Lab Causal Discovery Challenge organized by ADIA Lab and CrunchDAO~\cite{crunchdao_adia_docs, crunchdao_causal_2024}. The full competition dataset contains approximately 47,000 synthetic graph instances generated from structural causal models with nonlinear mechanisms including neural networks, Gaussian processes, and Bayesian models~\cite{olivetti2025adialab}. The dataset is divided into three partitions accessed through the CrunchDAO platform. The training set contains 23,500 dataset instances with provided ground-truth adjacency matrices, used for model training and local validation. The public test set contains 1,880 dataset instances scored by the CrunchDAO platform during the competition, providing an intermediate evaluation checkpoint before final submission. The private test set is a held-out partition revealed only at final submission; the proposed model achieves 81.082\% Balanced Accuracy on this partition, representing the headline result of this thesis.

% =====================================================================
\subsection{Label Derivation from the Adjacency Matrix}
\label{sec:label_derivation}
% =====================================================================

The eight-class label for each variable $K \notin \{X, Y\}$ is determined by the four binary entries of the ground-truth adjacency matrix $\mathbf{A}$ corresponding to edges incident to $K$ with respect to $X$ and $Y$: $\mathbf{A}_{K,X}$ (presence of edge $K \to X$), $\mathbf{A}_{X,K}$ (presence of edge $X \to K$), $\mathbf{A}_{K,Y}$, and $\mathbf{A}_{Y,K}$. Table~\ref{tab:label_rules_eval} enumerates the mapping from these four bits to the eight causal role labels. A Confounder satisfies $K \to X$ and $K \to Y$, a Mediator satisfies $X \to K$ and $K \to Y$, a Collider satisfies $X \to K$ and $Y \to K$, and the remaining four roles capture asymmetric relationships where $K$ involves exactly one of $X$ or $Y$ in a single directional edge. A detailed causal interpretation of each role is provided in Chapter~\ref{chap:background} (Section~\ref{sec:eight_roles}).

\begin{table}[ht]
	\centering
	\caption{Label derivation from the adjacency matrix for variable $K$. Each label is determined by the four binary entries of $\mathbf{A}$ corresponding to edges incident to $K$ with respect to $X$ and $Y$. The first row for each label shows the four bits in the order $K\to X,\; X\to K,\; K\to Y,\; Y\to K$. All variables in the ADIA Lab dataset are normalized to $[-1,1]$.}
	\label{tab:label_rules_eval}
	\begin{tabular}{lcccc}
		\toprule
		\textbf{Label} & $K\to X$ & $X\to K$ & $K\to Y$ & $Y\to K$ \\
		\midrule
		Confounder       & 1 & 0 & 1 & 0 \\
		Mediator         & 0 & 1 & 1 & 0 \\
		Collider         & 0 & 1 & 0 & 1 \\
		Cause of X       & 1 & 0 & 0 & 0 \\
		Cause of Y       & 0 & 0 & 1 & 0 \\
		Consequence of X & 0 & 1 & 0 & 0 \\
		Consequence of Y & 0 & 0 & 0 & 1 \\
		Independent      & 0 & 0 & 0 & 0 \\
		\bottomrule
	\end{tabular}
\end{table}

% =====================================================================
\subsection{Dataset Characteristics}
\label{sec:data_characteristics}
% =====================================================================

Each dataset instance contains $n = 1000$ observations over $p \in \{3,\ldots,10\}$ variables, with $X$ and $Y$ pre-identified as the treatment and outcome pair and the edge $X \to Y$ guaranteed to be present in the ground-truth DAG. Instances with $p > 3$ contribute more than one variable $K$ to the classification task, so the total number of $(K, \mathrm{label})$ training pairs exceeds 23,500. The eight classes are not uniformly distributed: Independent is the most frequent class while Mediator and Collider have the fewest training samples. This class imbalance motivates both the inverse-frequency class weights used in the loss function (Section~\ref{sec:loss}) and the use of Balanced Accuracy as the evaluation metric. The data are generated from six synthetic graph types and real-world food-web graphs, covering a range of causal structures that test the model's ability to distinguish fine-grained role differences~\cite{olivetti2025adialab}.

% =====================================================================
\section{Evaluation Metric: Balanced Accuracy}
\label{sec:metric}
% =====================================================================

The evaluation metric for all experiments is Balanced Accuracy (BA), defined as the mean per-class recall across all $C = 8$ classes:
\begin{equation}
\mathrm{BA} = \frac{1}{C} \sum_{c=1}^{C} \frac{\mathrm{TP}_c}{\mathrm{TP}_c + \mathrm{FN}_c}.
\label{eq:balanced_acc}
\end{equation}
Balanced Accuracy is the appropriate metric for three interconnected reasons. First, the class distribution is imbalanced with Independent as the majority class and Mediator and Collider as rare classes; reporting overall accuracy would overweight performance on the majority class and obscure failures on minority classes that are equally important for applications. Second, all eight causal roles carry equal weight in downstream decision-making -- misclassifying a Mediator as a Consequence of X has consequences for intervention planning that are no less serious than misclassifying an Independent as a Confounder. Third, an uninformed model that always predicts the most frequent class achieves approximately $\frac{1}{8} = 12.5\%$ Balanced Accuracy, providing a clear floor for the scale of improvement needed. The competition organizers' supervised baseline achieved approximately 58\% while the first-place solution achieved 76.70\%~\cite{olivetti2025adialab}.

% =====================================================================
\section{Experimental Baselines}
\label{sec:baselines}
% =====================================================================

% =====================================================================
\subsection{External Baselines}
\label{sec:external_baselines}
% =====================================================================

Two external baselines establish the competitive landscape. The competition first-place solution~\cite{adia_rank1_report} is the primary reference method, achieving 76.70\% Balanced Accuracy on the public leaderboard. Its architecture consists of 3-channel edge tensor construction, a stem Conv1D layer, five residual Conv1D blocks, global average pooling, edge-type embedding merge, two self-attention layers, and dual heads (8-class node classification plus binary edge auxiliary). All improvements in this thesis are measured against a faithful reimplementation achieving 73.96\% local and 73.61\% public leaderboard Balanced Accuracy, confirming reproducibility of the deep learning approach and providing a rigorous baseline for measuring incremental contributions.

The machine learning baselines represent the feature engineering and gradient boosting paradigm pushed to its limit. v10 applies LightGBM to eleven hand-crafted statistical features including Pearson correlation, mutual information, and ANM scores, achieving 63.11\%. v13 implements a comprehensive pipeline with over 300 engineered features combining correlation coefficients, partial correlations, conditional mutual information, HSIC scores, and outputs from four classical causal discovery algorithms (PC, LiNGAM, NOTEARS, GES), processed through an XGBoost/LightGBM/CatBoost ensemble with GNN refinement, achieving 72.64\%. These results demonstrate that scalar feature engineering -- regardless of the number of features or the sophistication of the ensemble -- cannot approach the performance of the deep learning approach that preserves the full distributional shape of conditional relationships.

% =====================================================================
\subsection{Internal Ablation Baselines}
\label{sec:internal_baselines}
% =====================================================================

The internal ablation baselines track incremental contributions through the development chain. The starting point is v2, the faithful 3-channel reimplementation achieving 73.96\% local and 73.61\% public. The 8-channel extension without structural attention (v8b) isolates the contribution of multi-bandwidth kernel regression and ANM residuals, achieving 76.94\%. Adding the structural attention bias (v11) raises this to 79.59\%, isolating the architectural contribution of topology-aware attention. Combining with XY augmentation (v11+xyaug) reaches 81.14\% local and 80.82\% public, serving as the previous best result. The 2D pipeline variants isolate visual representation contributions: v26c (2D node images only, no edge pipeline) demonstrates the necessity of edge context by collapsing to 51.51\%, while v26e base (full dual-pipeline without augmentation) achieves 80.47\%, confirming that the 2D visual features provide complementary information that the 1D edge pipeline alone cannot capture.

% =====================================================================
\section{Experimental Setup}
\label{sec:setup}
% =====================================================================

% =====================================================================
\subsection{Hardware Environment}
\label{sec:hardware}
% =====================================================================

All experiments are conducted on a server equipped with two NVIDIA RTX 5880 GPUs providing 48\,GB combined VRAM, a 48-core CPU for parallelizing data preprocessing across multiple workers, and 251\,GB of system RAM. Training uses Distributed Data Parallel across both GPUs to accelerate the training loop on the augmented dataset of approximately 263,000 samples. The model has approximately 280K parameters, comfortably fitting within the available GPU memory while allowing batch sizes of 64 for stable training.

% =====================================================================
\subsection{Two-Tier Evaluation Procedure}
\label{sec:evaluation_procedure}
% =====================================================================

The evaluation procedure follows a two-tier protocol with three distinct scoring mechanisms. The first tier is local validation: after each training run, the model is evaluated immediately on the 1,880 held-out instances from the public test split, which have known ground-truth adjacency matrices. This provides rapid feedback for architectural decisions without consuming submission quota on the CrunchDAO platform. The second tier is the CrunchDAO public leaderboard: after local validation confirms a positive result, selected model versions are submitted for official scoring, providing a fair comparison against other competition participants. The third mechanism is the private test submission: the final proposed model is submitted once to the private test partition, which remains hidden until the end of the competition, achieving 81.082\% as the definitive performance estimate. This two-tier protocol (local development followed by platform submission) prevents overfitting to the public leaderboard while ensuring reproducibility through the local validation split.

% =====================================================================
\subsection{Experiment Design Principles}
\label{sec:design_principles}
% =====================================================================

Experiments follow an ablation-first protocol to ensure every architectural decision is grounded in evidence. For each model variant, training and evaluation proceed on the base 23,500-sample dataset first. If the local result shows clear improvement over the corresponding ablation baseline, the experiment is repeated with XY augmentation (approximately 263K samples) before any public leaderboard submission. This ordering ensures that computational resources are allocated to configurations with demonstrated positive returns rather than being spent on speculative changes. The protocol also prevents overfitting to the public test set by limiting submissions to configurations that show unambiguous local improvement, a concern given that the public leaderboard represents only a fraction of the total test distribution.

% =====================================================================
\section{Complete Experiment Summary}
\label{sec:exp_summary}
% =====================================================================

Table~\ref{tab:all_results} presents all model versions with their local and public leaderboard scores, organized chronologically and grouped by experiment category. The table enables direct comparison of the contribution of each component: the gap between v2 (3-channel baseline) and v8b (8-channel) shows the effect of multi-bandwidth representation; the gap between v8b and v11 shows the effect of structural attention bias; the gap between v26e base and v26c shows the contribution of the edge context pipeline; and the gap between v26e base and v26e+xyaug shows the effect of X/Y remap augmentation. Detailed analysis of these results, including per-class breakdowns and mechanistic explanations, is presented in Chapter~\ref{chap:results}.

\begin{table}[ht]
	\centering
	\caption{Complete experiment results. Local = 1,880-sample test set; LB = public leaderboard; Private = hidden test (final only). $\Delta$v2 = change vs.\ v2 baseline (73.96\%).}
	\label{tab:all_results}
	\small
	\begin{tabularx}{\textwidth}{lXccc}
		\toprule
		\textbf{Version} & \textbf{Description} & \textbf{Local} & \textbf{LB} & \textbf{$\Delta$v2} \\
		\midrule
		\multicolumn{5}{l}{\textit{ML baselines}} \\
		v10 & LightGBM on 11 scalar features & 63.11\% & --- & $-10.85$ \\
		v13 & 300+ features, 4 algo outputs & 72.64\% & --- & $-1.32$ \\
		\midrule
		\multicolumn{5}{l}{\textit{DL reimplementation (3-channel baseline)}} \\
		v2 & Faithful 3ch reimplementation & 73.96\% & 73.61\% & --- \\
		\midrule
		\multicolumn{5}{l}{\textit{Multi-bandwidth extension (8-channel)}} \\
		v5m & +multi-bw kernel (5ch) & 75.44\% & 74.80\% & $+1.48$ \\
		v5m+xyaug & v5m + XY aug & 77.91\% & 77.55\% & $+3.95$ \\
		v8b & 8ch (5 kernel + 3 ANM) & 76.94\% & --- & $+2.98$ \\
		v8b+xyaug & v8b + XY aug & 80.26\% & 80.22\% & $+6.30$ \\
		\midrule
		\multicolumn{5}{l}{\textit{Structural attention bias}} \\
		v11 & v8b + struct.\ bias & 79.59\% & --- & $+5.63$ \\
		v11+xyaug & v11 + XY aug & 81.14\% & 80.82\% & $+7.18$ \\
		v11 ens & 3-seed ensemble & 80.92\% & 80.27\% & $+6.96$ \\
		\midrule
		\multicolumn{5}{l}{\textit{2D scatter density pipeline}} \\
		v26b & 8ch pairwise Conv2D & 79.17\% & --- & $+5.21$ \\
		v26c & 8ch 2D only (no edge) & 51.51\% & --- & $-22.45$ \\
		v26d & Edge loss favored ($\lambda=1.0$) & 80.37\% & --- & $+6.41$ \\
		v26d nf & Node loss favored ($\lambda=0.3$) & 78.77\% & --- & $+4.81$ \\
		v26e 12ch & 12ch (incl.\ coeff.\ ch.) & 80.42\% & --- & $+6.46$ \\
		v26e 8ch & 8ch (coeff.\ removed) & 80.47\% & --- & $+6.51$ \\
		\midrule
		\multicolumn{5}{l}{\textit{Final model (v26e 8ch + XY aug)}} \\
		v26e+xyaug & \textbf{Proposed} & \textbf{81.19\%} & \textbf{80.96\%} & $\mathbf{+7.23}$ \\
		v26e+xyaug & \textbf{Private} & --- & --- & \textbf{81.082\%} \\
		\bottomrule
	\end{tabularx}
\end{table}

% =====================================================================
\section*{Chapter Summary}
\addcontentsline{toc}{section}{Chapter Summary}
% =====================================================================

This chapter described the evaluation setup for assessing the proposed method. The dataset consists of 23,500 training instances and 1,880 local test instances from the ADIA Lab Causal Discovery Challenge, with a private test partition for final evaluation. Balanced Accuracy is the sole evaluation metric, defined as mean per-class recall across the eight causal roles, chosen because the class distribution is imbalanced and all eight roles carry equal decision-making weight. The baseline landscape spans two ML approaches (v10 at 63.11\% and v13 at 72.64\%), the reimplemented 3-channel deep learning baseline (v2 at 73.96\% local, 73.61\% public), the previous best configuration with structural attention bias (v11+xyaug at 81.14\% local, 80.82\% public), and the proposed final model (v26e 8ch + XY augmentation) achieving 81.082\% on the private test set. Detailed analysis of these results, including per-class breakdowns, ablation interpretation, and mechanistic explanations, is presented in Chapter~\ref{chap:results}.